\chapter{Abstract}

\begin{tabular}{@{} >{\bfseries}p{0.22\textwidth} p{0.75\textwidth} @{}}
\MakeUppercase{Keywords} &
Adaptive Protection, Inverter-Based Resources, Inverse Estimation,
Levenberg--Marquardt, Differential Protection, Overcurrent Protection,
Distance Protection, Microgrid Islanding, Digital Twin, Extended Kalman Filter,
Physics-Informed Neural Network, Wide-Area Protection Coordination, IEC~61850.
\end{tabular}
\par

\begin{doublespacing}

The global power grid is undergoing a fundamental structural transformation:
synchronous machines---whose high-inertia fault signatures underpinned a century
of deterministic relay design---are being displaced at accelerating pace by
Inverter-Based Resources (IBRs). Grid-Following (GFL) and Grid-Forming (GFM)
converters are software-controlled current sources whose fault currents are
limited to 1.1--2.0~pu by fast sub-cycle current limiters, and whose sequence
contributions are determined by firmware rather than physical machine reactance.
This paradigm shift systematically invalidates the foundational assumptions of
legacy overcurrent (ANSI~51/67), distance (ANSI~21), and differential (ANSI~87)
relay algorithms, creating a protection crisis across distribution and
sub-transmission networks.

This thesis proposes the \textbf{Self-Adaptive Model-Based Protection (SAMBP)}
architecture to address this crisis. The framework treats adaptive protection
as an \emph{inverse estimation problem}: each protection element maintains an
explicit physical model of its protected device and identifies current model
parameters from real-time measurements using the Levenberg--Marquardt algorithm.
Adaptation decisions are gated by a four-component confidence score and bounded
by a provably safe update law that preserves trip security under estimation
uncertainty. The framework is organised across four original contributions,
each validated with full numerical simulation results.

\textbf{Contribution~C1} (Chapter~\ref{ch:differential}) develops the
adaptive differential protection suite. The physics-constrained adaptive 87L
line differential relay modulates the bias slope $k_m^*(\kibr)$ in real time,
achieving TPR\,=\,1.000 and zero false trips across
$k_\mathrm{ibr} \in [0.03, 0.85]$~pu. The adaptive transformer differential
scheme replaces ad-hoc second-harmonic blocking with a Wald sequential
probability ratio test, providing formal false-trip and miss-rate guarantees
($\alpha = \beta = 0.001$).

\textbf{Contribution~C2} (Chapters~\ref{ch:oc_distance}--\ref{ch:microgrid})
develops adaptive overcurrent and distance protection. The SAMBP-SyncOC
framework applies Levenberg--Marquardt inverse estimation to synchronous
generator overcurrent protection, achieving $R^2 > 0.999$ and reducing
spurious relay pickups by 60\% with no security degradation. The
IBR-corrected quadrilateral distance characteristic achieves
$P_\mathrm{detect} = 99.3\%$ and $P_\mathrm{secure} = 99.97\%$. For
inverter-based microgrids, a ROCOF-supervised islanding detector achieves
NDZ\,$< 1\%$ and 114~ms isolation time, compliant with IEEE~1547-2018, with
GOOSE-triggered adaptive setting-group switching completing within 5~ms.

\textbf{Contribution~C3} (Chapter~\ref{ch:monitoring_ml}) delivers a
non-blocking Extended Kalman Filter digital twin that jointly estimates IBR
penetration ratio, line impedance, and fault location from IED phasor
measurements (RMSE\,$\leq 0.024$~pu, bounded by the Cramér--Rao Lower Bound),
without ever delaying a primary trip decision. An integrated physics-informed
neural network fault classifier embeds KVL/KCL constraints directly in the
training loss, achieving 96.1\% five-class accuracy with zero
physics-constraint violations, and adapts to fleet composition drift via
CUSUM-triggered online learning with $< 2$~pp catastrophic forgetting.

\textbf{Contribution~C4} (Chapter~\ref{ch:coordination_security}) integrates
all local algorithms into a Wide-Area Protection Coordinator (WAPC) that
provides PMU-based fault localisation (94.5\% correct-line identification),
automatic CTI-graded relay settings dispatch (100\% CTI-compliant in $< 3$~s),
and a closed-loop CUSUM$\to$settings$\to$GOOSE adaptive protocol completing
the full adaptation cycle in approximately 31~minutes. The cybersecurity
architecture augments GOOSE channels with HMAC-SHA256 authentication,
reducing attack risk by 54.9\% at 0.72~ms latency overhead, and establishes
the formal adversarial robustness framework for CT-injection attacks on the
PINN classifier.

The SAMBP framework is the first published architecture to integrate
model-based relay adaptation, digital-twin state estimation, physics-informed
classification, and closed-loop wide-area coordination across all protection
elements with provable safety guarantees.

\end{doublespacing}
